\chapter{INTRODUCCIÓN}

En la actualidad, muchos centros de salud de (San Marcos Carazo) enfrentan dificultades en la organización de la atención a pacientes, principalmente debido al manejo manual de turnos, largas filas y la falta de priorización adecuada de casos urgentes. Esta situación provoca retrasos, desorden y una atención menos eficiente.

Ante esta problemática, se propone el desarrollo de un Sistema de Gestión de Turnos para Centros de Salud de (San Marcos), el cual consiste en un prototipo de aplicación digital capaz de registrar pacientes, asignar turnos de manera automática y organizarlos según su nivel de prioridad.

Este prototipo representa una solución tecnológica inicial que contribuye a la modernización de los servicios de salud (San Marcos), promoviendo una atención más ordenada, rápida y eficiente para todos los usuarios.
